Overlapping sliding windows are a common localization device in persistent-homology-based change-point detection. Reusing the same raw observations across adjacent windows, however, alters the joint dependence structure of window-level features and can invalidate the exchangeability required by full permutation calibration. We first derive the exact covariance structure of overlapping window means and give a finite-sample counterexample showing that window-derived statistics may fail to be exchangeable even when the raw sequence is independent and identically distributed. We then compare naive permutation with a dependence multiplier bootstrap under two complementary budget regimes, fixed window count and fixed observation length, across 28 prespecified null configurations and two feature representations. At the nominal level 0.05, the per-configuration rejection rate of naive permutation reaches 0.9996. The multiplier procedure markedly reduces extreme anti-conservatism, yet its rejection rates still range from 0.0020 to 0.1412, leaving both conservative and anti-conservative regions. Under two opposite alternatives, a shape change and a mean shift, the size-standardized powers of the topological and moment representations exhibit an almost complete reversal. These findings characterize inferential validity under overlapping windows as a condition-dependent finite-sample calibration surface, and show that representation sensitivity and calibration validity should be evaluated separately while allowing for their finite-sample interaction.

\noindent\textbf{Keywords}: persistent homology; change-point detection; overlapping sliding windows; exchangeability; finite-sample calibration.
